\section{模型检验}

\subsection{结构与提交合同}
自动校验检查 8 份结果工作簿的工作表名、表头、合并单元格、列宽、两位小数、正式巷道编号和起终点坐标；
同时检查 Q4 从 $t=1$ 开始的完整前缀、路径时间单调性和出口可达性。全部 8 份工作簿和 12 条路线通过。

\subsection{守恒律与反例回归}
本题没有独立观测真值，因此采用守恒律、手算小图和单调性进行验证，结果见表\ref{tab:validation}。

\begin{table}[H]
\centering
\caption{模型不变量与反例回归验证}
\label{tab:validation}
\begin{tabular}{p{0.27\textwidth}p{0.53\textwidth}c}
\toprule
检验 & 预期与结果 & 状态\\
\midrule
对称菱形汇流 & 两支路各 $15$，汇流后恢复 $30\,\mathrm{m^3/min}$ & 通过\\
延迟同向水源 & 已湿节点新增水源后，下游流量升至 $60\,\mathrm{m^3/min}$ & 通过\\
相反端点水源 & 100 m 巷道双端供水的手算充满时刻为 21 min & 通过\\
虚拟拆边聚合 & 100 m 巷道在 10 m 处拆分，整巷充满时刻为 40 min & 通过\\
途中来水 & 1200 m 巷道第 1 min 来水，分段通行时间为 17 min & 通过\\
局部分段水向 & 中点水源两侧分别判为相反的顺水方向 & 通过\\
\bottomrule
\end{tabular}
\end{table}

全规模计算中，矿井 1 的 Q1/Q3 质量守恒绝对误差分别为
$7.92\times10^{-9}$、$1.30\times10^{-6}\,\mathrm{m^3}$，矿井 2 分别为
$8.73\times10^{-11}$、$1.16\times10^{-10}\,\mathrm{m^3}$；最大相对误差为
$1.76\times10^{-11}$。Q3 对所有 Q1 已到达端点和已充满巷道均未出现反常延后。
此外，23 项自动测试全部通过，覆盖投影吸附、质量守恒、相向水源、途中状态切换、Q4 连续复演和 Excel 合同。
